\documentclass[a4paper, serif, 10pt]{moderncv}

%% moderncv
% style and color
\moderncvstyle{banking}
\moderncvcolor{black}

%% page layout/margins
\usepackage[top=2.5cm, bottom=2.5cm, left=1.5cm, right=1.5cm]{geometry}

\nopagenumbers{}

%% Babel config for German language
% ref:
% - https://latex3.github.io/babel/guides/locale-german.html
\usepackage[ngerman]{babel}

%% fontspec
\usepackage{fontspec}

\IfFontExistsTF{Linux Libertine}{
  \setmainfont[Ligatures={TeX, Common}, Numbers=OldStyle]{Linux Libertine}
}{}
\IfFontExistsTF{Linux Libertine O}{
  \setmainfont[Ligatures={TeX, Common}, Numbers=OldStyle]{Linux Libertine O}
}{}

%% CTeX
% CJK support, used to show CJK characters in the resume
%
% - fontset=none: disable builtin fontset but instead set the CJK font manually
% - heading=false: disable ctex heading
% - punct=kaiming: use kaiming punctuations styles for CJK
% - scheme=plain: use plain scheme, do not override \normalsize font size
% - space=auto: space settings for CJK characters
%
% ref:
% - http://ctan.mirrorcatalogs.com/language/chinese/ctex/ctex.pdf
\usepackage[UTF8, heading=false, punct=kaiming, scheme=plain, space=auto]{ctex}

\IfFontExistsTF{Noto Serif CJK SC}{
  \setCJKmainfont{Noto Serif CJK SC}
}{}
\IfFontExistsTF{Noto Sans CJK SC}{
  \setCJKsansfont{Noto Sans CJK SC}
}{}

%% line spacing
\usepackage{setspace}
\setstretch{1.125}

%% URL styling - use normal text instead of monospace
\urlstyle{same}

% Auto-underline all links
% Defer until \begin{document} so hyperref is already loaded
\AtBeginDocument{%
  \let\oldhref\href
  \renewcommand{\href}[2]{\oldhref{#1}{\underline{#2}}}%
}

%% Patch \httplink and \httpslink to handle full URLs
% The original moderncv macros always prepend http:// or https:// to URLs.
% This causes issues when the URL already has a protocol (e.g., https://example.com)
% resulting in malformed URLs like https://https://example.com.
% This patch checks if the URL already contains "://" and uses it directly if so.
% Note: We use \str_set:Nx (x-expansion) to expand macros like \@homepage first.
\ExplSyntaxOn
\str_new:N \g_yamlresume_url_str

\RenewDocumentCommand{\httpslink}{O{}m}{
  \str_set:Nx \g_yamlresume_url_str {#2}
  \str_if_in:NnTF \g_yamlresume_url_str {://}
    { \url{#2} }
    { \url{https://#2} }
}

\RenewDocumentCommand{\httplink}{O{}m}{
  \str_set:Nx \g_yamlresume_url_str {#2}
  \str_if_in:NnTF \g_yamlresume_url_str {://}
    { \url{#2} }
    { \url{http://#2} }
}
\ExplSyntaxOff

\name{Lena Hoffmann}{}
\title{DevOps Engineer | Cloud-Infrastruktur \& Automatisierung}
\phone[mobile]{+49 30 98765432}
\email{lena.hoffmann@example.de}
\homepage{https://www.lena-hoffmann.dev}

\address{Musterstraße 12, Berlin, Berlin, Deutschland, 10115}{}{}

\extrainfo{{\small \faGithub}\ \href{https://github.com/lhoffmann-devops}{@lhoffmann-devops} {} {} {} • {} {} {} 
{\small \faLinkedin}\ \href{https://www.linkedin.com/in/lena-hoffmann-devops}{@lena-hoffmann-devops}}

\begin{document}

\maketitle

\section{Basisinformationen}

\cvline{}{\begin{itemize}
\item DevOps Engineer mit 7+ Jahren Erfahrung in der Automatisierung, Bereitstellung und dem Betrieb skalierbarer Cloud-Infrastrukturen.
\item Spezialisiert auf Kubernetes, Terraform, AWS und CI/CD-Pipelines; starke Kenntnisse in Infrastructure as Code und Monitoring.
\item Erfolgreiche Zusammenarbeit mit Entwicklungsteams zur Steigerung der Deployment-Frequenz und Reduzierung von Ausfallzeiten.
\item Analytisches Denken, Teamfähigkeit und Leidenschaft für Open-Source-Technologien.
\end{itemize}}
\section{Ausbildung}

\cventry{Okt. 2015–Sept. 2017}
        {Master, Informatik, Punkte: 1,5}
        {Technische Universität Berlin}
        {\url{https://www.tu.berlin}}
        {}
        {\begin{itemize}
\item Schwerpunkt auf verteilten Systemen und Cloud-Infrastrukturen
\item Masterarbeit über automatisierte Skalierung von Microservices in Kubernetes
\end{itemize}
\textbf{Kurse}: Verteilte Systeme, Cloud Computing, Betriebssysteme, Netzwerksicherheit, Datenbanktechnologie, Software-Architektur}

\cventry{Okt. 2012–Sept. 2015}
        {Bachelor, Angewandte Informatik, Punkte: 1,8}
        {Hochschule für Technik und Wirtschaft Berlin}
        {\url{https://www.htw-berlin.de}}
        {}
        {\begin{itemize}
\item Grundlagen in Softwareentwicklung und Systemadministration
\item Praktische Projekte zu Automatisierung und Scripting
\end{itemize}
\textbf{Kurse}: Programmierung, Algorithmen und Datenstrukturen, Rechnernetze, Linux-Systemadministration}
\section{Berufserfahrung}

\cventry{Apr. 2022–Heute}
        {Senior DevOps Engineer}
        {CloudOps GmbH}
        {\url{https://www.cloudops.example}}
        {}
        {\begin{itemize}
\item Konzeption und Betrieb einer Multi-Account-AWS-Landschaft mit Terraform und Terragrunt für über 40 Teams
\item Aufbau und Weiterentwicklung von GitOps-Pipelines mit Argo CD, Helm und Kubernetes zur automatisierten Bereitstellung
\item Einführung von OpenTelemetry und Prometheus/Grafana zur ganzheitlichen Observability, Reduzierung der MTTR um 35 \%
\item Mentoring von vier Junior-Engineers und Etablierung von Best Practices für Infrastructure as Code
\item Durchführung von Kostenoptimierungen durch Reserved Instances, Spot-Nutzung und Rechte-Zusammenfassung
\end{itemize}
\textbf{Schlüsselwörter}: AWS, Terraform, Kubernetes, GitOps, Observability}

\cventry{Juli 2019–März 2022}
        {DevOps Engineer}
        {DataForge AG}
        {\url{https://www.dataforge.example}}
        {}
        {\begin{itemize}
\item Automatisierung von Build-, Test- und Deployment-Prozessen mit Jenkins, GitLab CI und Docker
\item Migration von monolithischen Anwendungen in containerisierte Microservices auf Kubernetes
\item Implementierung von Security-Scans und Policy-as-Code mit OPA/Gatekeeper und Trivy
\item Aufbau von Infrastructure as Code mit Terraform und Ansible für AWS und On-Premises-Systeme
\item Sicherstellung der Einhaltung von SLAs durch proaktives Monitoring und Incident-Response
\end{itemize}
\textbf{Schlüsselwörter}: Jenkins, Docker, Kubernetes, Terraform, Ansible}

\cventry{Sept. 2017–Juni 2019}
        {Junior DevOps Engineer}
        {WebWerk Solutions}
        {\url{https://www.webwerk.example}}
        {}
        {\begin{itemize}
\item Betreuung von Linux-Servern (Debian, Ubuntu) und Automatisierung wiederkehrender Aufgaben mit Bash und Python
\item Aufbau zentraler Logging-Lösung mit ELK Stack für über 50 Anwendungen
\item Unterstützung bei der Einführung von CI/CD mit GitLab und automatisierten Tests
\item Erstellung von Monitoring-Dashboards und Alarmen mit Nagios und Icinga
\end{itemize}
\textbf{Schlüsselwörter}: Linux, Bash, Python, ELK, CI/CD}
\section{Sprachen}

\cvline{Deutsch}{Muttersprache oder zweisprachig \hfill \textbf{Schlüsselwörter}: Muttersprache}
\cvline{Englisch}{Verhandlungssichere Kenntnisse \hfill \textbf{Schlüsselwörter}: TOEFL iBT 110, Technisches Englisch}
\cvline{Französisch}{Gute Kenntnisse \hfill \textbf{Schlüsselwörter}: Grundkenntnisse}
\section{Fähigkeiten}

\cvline{Cloud \& Infrastruktur}{Experte \hfill \textbf{Schlüsselwörter}: AWS, Azure, Terraform, CloudFormation, Pulumi}
\cvline{Container \& Orchestrierung}{Experte \hfill \textbf{Schlüsselwörter}: Kubernetes, Docker, Helm, Argo CD, Istio}
\cvline{CI/CD \& Automatisierung}{Erfahren \hfill \textbf{Schlüsselwörter}: GitLab CI, Jenkins, GitHub Actions, Ansible, Argo Workflows}
\cvline{Monitoring \& Observability}{Erfahren \hfill \textbf{Schlüsselwörter}: Prometheus, Grafana, OpenTelemetry, ELK Stack, Datadog}
\cvline{Programmierung \& Scripting}{Erfahren \hfill \textbf{Schlüsselwörter}: Python, Go, Bash, YAML, SQL}
\section{Auszeichnungen}

\cventry{Dez. 2023}
        {CloudOps GmbH}
        {DevOps Excellence Award}
        {}
        {}
        {Auszeichnung für die herausragende Automatisierung der Cloud-Infrastruktur und die Reduzierung der Betriebskosten um 25 \%.}

\cventry{Okt. 2017}
        {Technische Universität Berlin}
        {Beste Masterarbeit}
        {}
        {}
        {Prämierung der Masterarbeit zur automatisierten Skalierung von Microservices in Kubernetes.}
\section{Zertifikate}

\cventry{Mai 2023}
        {Amazon Web Services}
        {AWS Certified DevOps Engineer – Professional}
        {\url{https://aws.amazon.com/certification/}}
        {}
        {}

\cventry{Feb. 2022}
        {Cloud Native Computing Foundation}
        {Certified Kubernetes Administrator (CKA)}
        {\url{https://www.cncf.io/certification/cka/}}
        {}
        {}

\cventry{Nov. 2021}
        {HashiCorp}
        {HashiCorp Certified: Terraform Associate}
        {\url{https://www.hashicorp.com/certification/terraform-associate}}
        {}
        {}
\section{Veröffentlichungen}

\cventry{Sept. 2023}
        {Effiziente Kubernetes-Skalierung durch adaptive Metriken}
        {DevOps-Konferenz Berlin 2023}
        {\url{https://www.devops-konferenz-berlin.example/vortraege}}
        {}
        {\begin{itemize}
\item Vortrag über die Nutzung von Custom Metrics und KEDA zur automatisierten Skalierung von Workloads
\item Praxisbericht zur Reduzierung von Overprovisioning und Steigerung der Ressourceneffizienz
\end{itemize}}
\section{Projekte}

\cventry{Jan. 2023–Dez. 2023}
        {Eine interne Plattform zur Bereitstellung standardisierter Cloud-Umgebungen via Terraform und Kubernetes.}
        {Self-Service-Infrastrukturplattform}
        {\url{https://github.com/lhoffmann-devops/self-service-infra}}
        {}
        {\begin{itemize}
\item Entwicklung eines Backstage-Plugins für Self-Service-Deployments
\item Integration von Terraform-Modulen und GitOps-Workflows
\item Reduzierung der Bereitstellungszeit neuer Umgebungen von Tagen auf unter 30 Minuten
\end{itemize}
\textbf{Schlüsselwörter}: Terraform, Backstage, Kubernetes, GitOps}

\cventry{März 2022–Aug. 2022}
        {Benchmark und Optimierung des Kubernetes Cluster-Autoscalers für gemischte Workloads.}
        {Cluster-Autoscaler-Optimierung}
        {\url{https://github.com/lhoffmann-devops/cluster-autoscaler-benchmark}}
        {}
        {\begin{itemize}
\item Analyse von Skalierungsverhalten unter Last und Implementierung von PriorityClasses
\item Senkung der AWS-Kosten um 18 \% bei gleichbleibender Performance
\end{itemize}
\textbf{Schlüsselwörter}: Kubernetes, AWS, Autoscaling, Kostenoptimierung}
\section{Interessen}

\cvline{Open Source}{Kubernetes, Terraform, Prometheus}
\cvline{Sport}{Laufen, Klettern, Radfahren}
\cvline{Musik}{Klavier, Konzerte}
\section{Ehrenamtliche Tätigkeiten}

\cventry{Jan. 2021–Dez. 2023}
        {Mentorin für DevOps}
        {Open Source Helpdesk Berlin}
        {\url{https://www.oss-helpdesk-berlin.example}}
        {}
        {\begin{itemize}
\item Unterstützung von Einsteigern bei der Einrichtung von CI/CD-Pipelines und Kubernetes-Clustern
\item Organisation monatlicher Workshops zu Infrastructure as Code und Cloud-Sicherheit
\item Beratung von gemeinnützigen Organisationen bei der Migration in die Cloud
\end{itemize}}
    
\end{document}